\documentclass[9pt,twocolumn]{article}
\usepackage[letterpaper,margin=0.55in,top=0.48in,bottom=0.52in]{geometry}
\usepackage[T1]{fontenc}\usepackage{newtxtext,newtxmath}\usepackage{microtype}
\usepackage{amsmath,mathtools}\usepackage{booktabs,tabularx,array}\usepackage[normalem]{ulem}\usepackage{cancel}\usepackage{enumitem,xcolor,ragged2e}
\usepackage{graphicx}
\setlength{\columnsep}{0.24in}\setlength{\parindent}{0pt}\setlength{\parskip}{2.2pt}
\setlist{nosep,leftmargin=*}\renewcommand{\arraystretch}{1.10}
\newcommand{\oldclaim}[1]{\sout{#1}}
\newcommand{\claimbox}[1]{\begin{center}\fbox{\parbox{0.94\columnwidth}{\centering\scriptsize #1}}\end{center}}
\newcommand{\eclass}[1]{\textsf{[#1]}}
\DeclareFontFamilySubstitution{TS1}{ntxtlf}{cmr}
% \let binds the CURRENT meaning of \_ ; \newcommand would create a macro that
% re-expands to \_ and loops until TeX runs out of input stack.
\let\CIGunderscore\_
\renewcommand{\_}{\CIGunderscore\allowbreak}
\begin{document}
\sloppy
\twocolumn[{\begin{center}
{\LARGE\bfseries Report 4 - Structural Rankability}\\[2pt]
{\large\itshape Append-only update D: the family carrying the lower bound has no tolerance trade-off at all, and curve richness tracks coefficient precision rather than width}\\[2pt]
{\small CIG-AMF Scientific Evolution Report -- 11 September}
\end{center}\vspace{2mm}}]

\textbf{Append-only rule.} Sections 1--71, including addenda A--C, are an immutable prefix. This addendum begins at Section 72.

\section*{72. The staircase family has exactly two levels}
The programme's central object is $\chi_\varepsilon$, the number of rankings needed to stay within $\varepsilon$ of the budget-wise optimum after normalising the total component radius. The exact curve was computed for the staircase family and for disjoint unions of staircase blocks, over $\mathbb Q$, by testing directly whether a budget set is coverable by one chain and then solving the resulting set cover exactly.
\begin{center}\scriptsize
\begin{tabularx}{0.98\columnwidth}{X r r}
\toprule
blocks & $\chi$ below threshold & threshold\\
\midrule
$[2]$ / $[3]$ / $[4]$ & $2$ / $3$ / $4$ & $1/4$, $1/6$, $1/8$\\
$[2,2]$ / $[3,2]$ & $3$ / $3$ & $1/8$, $1/10$\\
$[3,3]$ / $[4,2]$ & $4$ / $4$ & $1/12$, $1/12$\\
\bottomrule
\end{tabularx}
\end{center}
Every configuration has \textbf{exactly two levels}, and the threshold is always precisely $1/(2\sum_b r_b)$. No intermediate value of $\chi$ occurs at any $\varepsilon$.
\claimbox{\textbf{This is a collapse, not a trade-off.} The family that carries the lower bound exhibits no $K$-versus-$\varepsilon$ curve: exact optimality needs $r$ rankings, and a tolerance of $1/(2r)$ --- already tiny --- needs one. Disjoint unions do not help; $[2,2]$ gives $3$, not $4$.}

\section*{73. The proved lower bound is loose by exactly a factor two}
\texttt{normalized\_lower} establishes $\chi_\varepsilon\ge r$ for $\varepsilon\le1/(4r)$. The exact collapse point is $1/(2r)$, so $\chi_\varepsilon=r$ holds on the whole interval $\varepsilon<1/(2r)$. The proved statement can be strengthened by a factor two, and this is the cheapest available Lean target in the chain. The matching upper bound $\lfloor2/\varepsilon\rfloor+1$ is meanwhile vacuous on this family: at the threshold it returns $4r+1$ where the truth is $1$.

\section*{74. Curve richness is governed by coefficient precision}
The scalar two-support-point model was swept with random integer weights of $B$ bits, $30$ draws per cell, counting levels of the exact $\chi_\varepsilon$ curve. The model reproduces the \texttt{ScalarChi5} witness to the fraction ($5\to3\to2\to1$ at $0$, $1/114$, $1/57$, $7/57$), so the measurement is calibrated.
\begin{center}\scriptsize
\begin{tabularx}{0.98\columnwidth}{r r r r r r r}
\toprule
$m$ & $B{=}1$ & $2$ & $3$ & $4$ & $5$ & $6$\\
\midrule
$5$ & $\mathbf{2}$ & $3$ & $3$ & $3$ & $3$ & $3$\\
$6$ & $\mathbf{2}$ & $3$ & $3$ & $3$ & $3$ & $3$\\
$7$ & $\mathbf{2}$ & $3$ & $3$ & $3$ & $\mathbf{4}$ & $\mathbf{4}$\\
\bottomrule
\end{tabularx}
\end{center}
At one bit the level count is capped at two for every width tested. Width contributes weakly; precision contributes sharply.
\claimbox{\textbf{The $\chi\ge5$ witness is not a geometric outlier.} Its four levels are an extreme draw from the random distribution at the same $m$ and the same bit width. This is exactly the failure mode the 7 September plan warned of: a witness whose difficulty may come from cleverly chosen coefficients rather than from support geometry.}

\section*{75. Consequence: the question has to change}
The claim ``how many rankings are forced, as a function of $m$'' now rests on two weak legs --- a family with a law that is trivial (two levels, threshold $1/m$), and a single witness whose richness traces to precision. The quantity the data actually points at is different:
\[\text{levels of }\chi_\varepsilon\ \longleftrightarrow\ \text{bit precision of the coefficients, at fixed dimension and normalisation.}\]
That is still a law, and still about representability, but it is not the law being sought. Stating it correctly would also change what a counterexample looks like: a family with unbounded level count at bounded precision would be the real result, and none of the families enumerated here is one.
\eclass{EXACT RATIONAL COMPUTATION ON NAMED FAMILIES} --- \texttt{research/structural\_tolerance\_20260910/tolerance\_curves.json}. No claim is made about families outside those enumerated.
\end{document}
